% Review draft; insert after the manuscript's existing \appendix command.
% Requires the companion main-text snippet and existing amsmath/amsthm setup.
% Each prose sentence occupies a separate source line.
\section{Derivations and supplementary details for the algorithmic enhancements}
\label{app:enhancements}

\subsection{Extended polymatroid inequalities}
\label{app:enh-queue}

The following derivation assumes the fixed task rates used in the planning model and the complete physical path domain, including all admissible return bases.
Write $f=f^{\mathrm{fleet}}>0$ within this subsection.
Let $n_{rq}$ denote the number of occurrences of OD pair $r$ in task $q$.
Demand conservation gives $h_r=\sum_q n_{rq}\hbar_q$, and $g_q=\sum_r n_{rq}c_r^{\mathrm{direct}}$.
In particular, a task containing two orders from the same OD pair contributes twice to both expressions.

\begin{proposition}
\label{prop:enh-queue-bound}
For every integer design $\boldsymbol d\in\mathcal X$, inequality~\eqref{eq:enh-queue-bound} holds.
Moreover, each function $v_l$ in \eqref{eq:enh-queue-function} is normalized and submodular.
\end{proposition}

\begin{proof}
Fix an integer design $\boldsymbol d$ and any feasible operating solution $u$, and denote its cost by $c_{\mathrm{op}}(\boldsymbol d,u)$.
Eliminating direct-courier flows through demand conservation and using the fleet requirement gives
\begin{equation}
    \begin{aligned}
    c_{\mathrm{op}}(\boldsymbol d,u)
    \ge {}& C^{\mathrm R}+f\sum_{l\in\mathcal L}I_l\\
    &+\sum_{q\in\hat{\mathcal R}}\sum_{l\in\mathcal L}
      \sum_{p\in\mathcal P_{ql}}
      \bigl(c_p+f\bar T_p-g_q\bigr)\psi_p\\
    &+\sum_{k\in\mathcal K}(c_k+f\bar T_k)\varphi_k.
    \end{aligned}
    \label{eq:enh-app-cost}
\end{equation}
For a nonempty task path set, its contribution to the second line is at least $-S_{ql}\alpha_{ql}\rho_l$, by \eqref{eq:enh-task-savings} and the task-flow equality $\sum_{p\in\mathcal P_{ql}}\psi_p=\hbar_q\alpha_{ql}\rho_l$.
If $\hbar_q=0$, all associated service flows vanish.
If $\hbar_q>0$ and $\mathcal P_{ql}=\varnothing$, the same equality forces $\alpha_{ql}\rho_l=0$.
Hence, the bound also holds under the stated zero-coefficient conventions.
Using $I_l\ge\rho_l/(1-\rho_l)$ and dropping the nonnegative repositioning contribution yields
\begin{equation}
    c_{\mathrm{op}}(\boldsymbol d,u)
    \ge C^{\mathrm R}+\sum_{l\in\mathcal L}
    \left[f\frac{\rho_l}{1-\rho_l}
    -\rho_l\sum_q S_{ql}\alpha_{ql}\right].
    \label{eq:enh-app-rho-bound}
\end{equation}
For $s\ge0$, define
\begin{equation}
    \Phi(s)
    :=\sup_{0\le\rho<1}\left\{s\rho-f\frac{\rho}{1-\rho}\right\}
    =\begin{cases}
      0,&0\le s\le f,\\
      (\sqrt{s}-\sqrt f)^2,&s>f.
    \end{cases}
    \label{eq:enh-app-phi}
\end{equation}
The maximizer is $\rho=0$ for $s\le f$ and $\rho=1-\sqrt{f/s}$ otherwise.
Minimizing the right-hand side of \eqref{eq:enh-app-rho-bound} independently over $0\le\rho_l<1$, and then taking the infimum over feasible operating solutions on the left, proves \eqref{eq:enh-queue-bound}.
Any additional restriction on availability in the original model can only increase the minimum operating cost.

The function $\Phi$ is convex because it is the pointwise supremum of affine functions of $s$.
Therefore, for $w\ge0$, the increment $\Phi(s+w)-\Phi(s)$ is nondecreasing in $s$.
Since $v_l(A)=-\Phi(\sum_{q\in A}S_{ql})$ and $S_{ql}\ge0$, for $A\subseteq B$ and $q\notin B$ we obtain
\begin{equation}
    v_l(A\cup\{q\})-v_l(A)
    \ge v_l(B\cup\{q\})-v_l(B).
    \label{eq:enh-app-submodular}
\end{equation}
This establishes submodularity, while $\Phi(0)=0$ gives $v_l(\varnothing)=0$.
\end{proof}

\begin{proposition}
\label{prop:enh-epi}
Inequality~\eqref{eq:enh-queue-epi} is valid for the integer operating-cost epigraph and its convex hull for every choice of the permutations $\pi_l$.
For fixed coefficients $S_{ql}$, a most violated inequality in this family is obtained by sorting the current assignment values in nonincreasing order for each base.
\end{proposition}

\begin{proof}
Fix a base and suppress its index.
For a permutation $\pi$, let $A_j$ be its prefix sets and $m_{\pi(j)}=v(A_j)-v(A_{j-1})$.
For any task subset $T$ and each $\pi(j)\in T$, submodularity implies
\[
    v(T\cap A_j)-v(T\cap A_{j-1})
    \ge v(A_j)-v(A_{j-1})=m_{\pi(j)}.
\]
Summing these inequalities and using $v(\varnothing)=0$ gives $\sum_{q\in T}m_q\le v(T)$.
Applying this result to $T=A_l(\boldsymbol d)$ at each base and using \eqref{eq:enh-queue-bound} proves validity at every integer epigraph point.
Linearity extends validity to its convex hull.

To prove exact separation, set $x_q=\bar\alpha_{ql}$ and choose $\pi$ such that $x_{\pi(1)}\ge\cdots\ge x_{\pi(n)}\ge0$, with $x_{\pi(n+1)}=0$.
Let $a$ be the marginal coefficient vector associated with any other permutation, and write $a(A)=\sum_{q\in A}a_q$.
The preceding argument gives $a(A_j)\le v(A_j)$, so summation by parts yields
\begin{equation}
    \begin{aligned}
    a^\top x
    &=\sum_{j=1}^{n}(x_{\pi(j)}-x_{\pi(j+1)})a(A_j)\\
    &\le\sum_{j=1}^{n}(x_{\pi(j)}-x_{\pi(j+1)})v(A_j)
    =m^\top x.
    \end{aligned}
    \label{eq:enh-app-separation}
\end{equation}
The last equality follows from $m(A_j)=v(A_j)$ for the sorted permutation.
The permutations can be selected independently across bases, establishing the claim for \eqref{eq:enh-queue-epi}.
Ties may be broken arbitrarily.
Once the coefficients $S_{ql}$ are available, separation requires $O(|\mathcal L|n\log n)$ time.
\end{proof}

The exact-separation statement concerns the permutation inequalities for the fixed functions $v_l$.
It does not assert submodularity of the full operating value function or a convex-hull description of the complete planning model.

\subsection{Disjunctive cuts}
\label{app:enh-state}

For a fixed task $q$, the state sets used in the main text are
\begin{equation}
    \begin{aligned}
    \mathcal X_{q0}
      &=\{\boldsymbol d\in\mathcal X:\alpha_{ql}=0\ \forall l\in\mathcal L\},\\
    \mathcal X_{ql}
      &=\{\boldsymbol d\in\mathcal X:\alpha_{ql}=1,\
         \alpha_{qm}=0\ \forall m\in\mathcal L\setminus\{l\}\},
         \quad l\in\mathcal L.
    \end{aligned}
    \label{eq:enh-app-states}
\end{equation}
All other design variables remain subject to their original constraints and are unfixed.
For each state, the required total-cost bound satisfies
\begin{equation}
    b_{qs}\le
    \inf_{\boldsymbol d\in\mathcal X_{qs}}
    \left\{F(\boldsymbol d)+Q(\boldsymbol d)\right\},
    \qquad s\in\{0\}\cup\mathcal L.
    \label{eq:enh-state-bound}
\end{equation}

\begin{proposition}
\label{prop:enh-state}
For any finite bounds satisfying \eqref{eq:enh-state-bound}, inequality~\eqref{eq:enh-state-cut} is valid for the integer total-cost epigraph and its convex hull.
\end{proposition}

\begin{proof}
Let $u_{q0}=1-\sum_{l\in\mathcal L}\alpha_{ql}$ and $u_{ql}=\alpha_{ql}$ for $l\in\mathcal L$.
At every integer design, these indicators are nonnegative and sum to one, with exactly one equal to one.
If that indicator corresponds to state $s$, then
\[
    F(\boldsymbol d)+\zeta
    \ge F(\boldsymbol d)+Q(\boldsymbol d)
    \ge b_{qs}
    =\sum_{t\in\{0\}\cup\mathcal L}b_{qt}u_{qt}.
\]
Substituting the state indicators gives \eqref{eq:enh-state-cut}, and linearity establishes validity for the convex hull.
\end{proof}

For fixed state constants $b_{qs}$, the inequality also describes the lower boundary of the convex hull of the state epigraphs $\{(u,T):u=e_s,\ T\ge b_{qs}\}$, where $T$ denotes total cost.
Indeed, a convex combination with weights $u_s$ has minimum ordinate $\sum_s b_{qs}u_s$.
This statement concerns the relaxation defined by the state constants, rather than the full planning model.

To obtain $b_{qs}$, we fix the assignment state and relax the remaining discrete design variables while retaining the complete physical operating domain.
A certified lower bound for this relaxation satisfies \eqref{eq:enh-state-bound}.
When the relaxation is solved by column generation, certification requires pricing over all admissible paths or a valid correction for omitted columns.
The objective of a restricted-column feasible solution is insufficient for this purpose.
An existing certified global total-cost lower bound provides a valid fallback for any state whose stronger bound has not been certified.
Each selected task generates a separate inequality on the same total-cost epigraph.

\subsection{Column reuse and incumbent initialization}
\label{app:enh-warm}

Generated service and repositioning paths are retained in a shared column pool and supplied as initial columns for subsequent restricted subproblems on the same physical path domain.
Each subproblem is reoptimized for its current design, and complete-domain pricing remains responsible for lower-bound certification.
The best available feasible design and its operating cost are supplied to the master as a MIP start.

For the two-base formulation with single- and two-order tasks, an auxiliary adaptive large neighborhood search constructs and improves these incumbents.
Neighborhoods are repaired by solving restrictions of the original MISOCP to selected path columns, retaining the physical flow, fleet, queueing, and capacity constraints.
Accepted candidates are checked for feasibility and evaluated under the original objective before updating the incumbent.
These restricted repairs provide upper-bound candidates; their local bounds are not used as global lower bounds.
